\section{Validation}
\label{sec:validation}

\subsection{Li et al.\ (2023) Benchmark Cases}
Validation uses experimental data from Li et al.\ \cite{li2023} Table~2 (Q345 steel; Table~1 properties in \cref{sec:methods}). The inherent-strain amplitude $\varepsilon_0$ is calibrated on four anchor cases (1, 6, 87, 90) and applied without re-fitting to all published Table~2 rows available in our dataset: cases 1--7 and 87--90 (rows 8--86 are omitted in the source excerpt). We therefore report calibration-anchor, non-anchor validation, and all-case metrics separately.

\Cref{tab:li2023-exp} lists experimental conditions. \Cref{tab:li2023-all} reports predicted vs.\ experimental midspan residual deflection for all eleven simulations in the moving-torch baseline. \Cref{tab:li2023-cal} highlights the calibration subset. A later process comparison, \cref{tab:li2023-process}, evaluates a full-line heat-source approximation with quenching.

% --- Flowchart styling (shared with publication/sections/flowcharts.tex) ---
\tikzset{
  block/.style={
    rectangle,
    rounded corners=2pt,
    draw=black!70,
    fill=blue!7,
    align=center,
    minimum height=9mm,
    text width=36mm,
    font=\small
  },
  wideblock/.style={
    block,
    text width=52mm
  },
  smallblock/.style={
    block,
    text width=29mm,
    font=\footnotesize
  },
  decision/.style={
    diamond,
    aspect=2.1,
    draw=black!70,
    fill=orange!12,
    align=center,
    inner sep=1.5pt,
    text width=26mm,
    font=\small
  },
  group/.style={
    rectangle,
    rounded corners=3pt,
    draw=black!40,
    fill=gray!5,
    inner sep=7pt
  },
  arrow/.style={
    -{Latex[length=2.4mm]},
    thick,
    draw=black!75
  },
  dashedarrow/.style={
    arrow,
    dashed,
    draw=black!55
  },
  note/.style={
    rectangle,
    rounded corners=2pt,
    draw=black!45,
    fill=yellow!10,
    align=left,
    text width=130mm,
    font=\small
  }
}

\begin{figure}[H]
\centering
\resizebox{\linewidth}{!}{%
\begin{tikzpicture}[node distance=11mm and 13mm]
  \node[wideblock] (data) {Li et al.\ 2023\\Table~1 material + Table~2 cases};

  \node[wideblock, below left=17mm and 10mm of data] (v9) {v9 torch baseline\\moving Gaussian\\sequential passes};
  \node[wideblock, below=17mm of data] (v10) {v10 full-line flux\\collapsed multi-pass + quench};
  \node[wideblock, below right=17mm and 10mm of data] (v11) {v11 full-line flux\\sequential passes + quench};

  \node[wideblock, below=of v9] (r9) {All 11 Li2023 cases\\about 1.2\% MAPE\\primary baseline};
  \node[wideblock, below=of v10] (r10) {Multi-pass cases\\about $-$50\% bias\\schedule collapsed};
  \node[wideblock, below=of v11] (r11) {Cases 6, 7, 87--89\\about 0.3\% MAPE\\schedule preserved};

  \node[wideblock, below=20mm of r10] (paper) {Validation metrics, figures, discussion\\manuscript};

  \draw[arrow] (data) -- (v9);
  \draw[arrow] (data) -- (v10);
  \draw[arrow] (data) -- (v11);
  \draw[arrow] (v9) -- (r9);
  \draw[arrow] (v10) -- (r10);
  \draw[arrow] (v11) -- (r11);
  \draw[arrow] (r9) -- (paper);
  \draw[arrow] (r10) -- (paper);
  \draw[arrow] (r11) -- (paper);
\end{tikzpicture}%
}
\caption{Validation workflow for Li et al.\ (2023) cases and the full-line thermal-flux process comparison variants.}
\label{fig:li2023-validation-flow}
\end{figure}

\begin{table}[H]
\centering
\caption{Li et al.\ (2023) Table~2: published experimental cases used in this study.}
\label{tab:li2023-exp}
\small
\begin{tabular}{rrrrrr}
\toprule
Case & $v$ (mm/s) & $E$ (J/mm) & $t$ (mm) & Passes & $w_{\mathrm{exp}}$ (mm) \\
\midrule
1  & 6  & 600  & 14 & 1 & 5.785 \\
2  & 7  & 600  & 14 & 1 & 5.771 \\
3  & 8  & 600  & 14 & 1 & 5.760 \\
4  & 9  & 600  & 14 & 1 & 5.748 \\
5  & 10 & 600  & 14 & 1 & 5.727 \\
6  & 11 & 600  & 14 & 2 & 7.712 \\
7  & 12 & 600  & 14 & 2 & 7.696 \\
\multicolumn{6}{c}{$\vdots$ (cases 8--86 omitted in source)} \\
87 & 12 & 1000 & 10 & 2 & 13.506 \\
88 & 13 & 1000 & 10 & 2 & 13.443 \\
89 & 14 & 1000 & 10 & 2 & 13.383 \\
90 & 15 & 1000 & 10 & 1 & 11.322 \\
\bottomrule
\end{tabular}
\end{table}

\begin{table}[H]
\centering
\caption{Full Li2023 validation: predicted vs.\ experimental deflection (all eleven published cases).}
\label{tab:li2023-all}
\footnotesize
\resizebox{\linewidth}{!}{%
\begin{tabular}{rrrrrrrrr}
\toprule
Case & $v$ & $E$ & $t$ & $w_{\mathrm{exp}}$ & $w_{\mathrm{pred}}$ & $\Delta w$ & $\Delta w/w_{\mathrm{exp}}$ (\%) & $T_{\mathrm{peak}}$ (K) \\
 & (mm/s) & (J/mm) & (mm) & (mm) & (mm) & (mm) &  &  \\
\midrule
1  & 6  & 600  & 14 & 5.785 & 5.888 &  0.103 &  1.77 & 1174 \\
2  & 7  & 600  & 14 & 5.771 & 5.888 &  0.117 &  2.02 & 1174 \\
3  & 8  & 600  & 14 & 5.760 & 5.887 &  0.127 &  2.21 & 1189 \\
4  & 9  & 600  & 14 & 5.748 & 5.886 &  0.138 &  2.40 & 1174 \\
5  & 10 & 600  & 14 & 5.727 & 5.886 &  0.159 &  2.78 & 1173 \\
6  & 11 & 600  & 14 & 7.712 & 7.722 &  0.010 &  0.13 & 1177 \\
7  & 12 & 600  & 14 & 7.696 & 7.722 &  0.026 &  0.34 & 1175 \\
87 & 12 & 1000 & 10 & 13.506 & 13.486 & $-$0.020 & $-$0.15 & 1388 \\
88 & 13 & 1000 & 10 & 13.443 & 13.486 &  0.043 &  0.32 & 1384 \\
89 & 14 & 1000 & 10 & 13.383 & 13.485 &  0.102 &  0.76 & 1386 \\
90 & 15 & 1000 & 10 & 11.322 & 11.319 & $-$0.003 & $-$0.03 & 1376 \\
\bottomrule
\end{tabular}
}
\end{table}

\noindent Across all eleven cases: mean $|\Delta w|=\SI{0.077}{mm}$, RMSE $=\SI{0.094}{mm}$, bias $=\SI{0.073}{mm}$, MAPE $=\SI{1.17}{\%}$, and maximum $|\Delta w/w_{\mathrm{exp}}|=\SI{2.78}{\%}$ (Case~5). Multi-pass and high-energy cases (6, 7, 87--90) agree within \SI{0.8}{\%}. Single-pass speed-sweep cases (1--5) show a systematic overprediction and weak sensitivity to $v$ in $w_{\mathrm{pred}}$ (discussed in \cref{sec:discussion}).

\begin{table}[H]
\centering
\caption{Calibration subset (cases 1, 6, 87, 90) used to tune $\varepsilon_0$.}
\label{tab:li2023-cal}
\begin{tabular}{rrrrrrrr}
\toprule
Case & $w_{\mathrm{exp}}$ (mm) & $w_{\mathrm{pred}}$ (mm) & $\Delta w$ (mm) & $\Delta w/w_{\mathrm{exp}}$ (\%) & $\varepsilon_0$ \\
\midrule
1  & 5.785 & 5.888 &  0.103 &  1.77 & 0.0720 \\
6  & 7.712 & 7.722 &  0.010 &  0.13 & 0.0477 \\
87 & 13.506 & 13.486 & $-$0.020 & $-$0.15 & 0.0582 \\
90 & 11.322 & 11.319 & $-$0.003 & $-$0.03 & 0.0978 \\
\bottomrule
\end{tabular}
\end{table}

\noindent Calibration subset: mean $|\Delta w|=\SI{0.034}{mm}$, MAPE $=\SI{0.52}{\%}$. Non-anchor validation cases (2--5, 7, 88, 89) have mean $|\Delta w|=\SI{0.102}{mm}$, RMSE $=\SI{0.112}{mm}$, bias $=\SI{0.102}{mm}$, and MAPE $=\SI{1.55}{\%}$.

\begin{table}[H]
\centering
\caption{Mesh/time-step sensitivity study for representative Li2023 Case~6 ($v=\SI{11}{mm/s}$, $E=\SI{600}{J/mm}$, $t=\SI{14}{mm}$, two passes). Baseline mesh uses $h=\SI{20}{mm}$, $h_{\mathrm{ref}}=\SI{5}{mm}$, and $b_{\mathrm{ref}}=\SI{60}{mm}$; refined mesh uses $h_{\mathrm{ref}}=\SI{2.5}{mm}$ with the same background spacing and band width.}
\label{tab:li2023-mesh-dt-sensitivity}
\small
\begin{tabular}{lrrr}
\toprule
Setting & $\Delta t$ (s) & $T_{\mathrm{peak}}$ (K) & $|w|_{\max}$ (mm) \\
\midrule
Baseline (paper setting) & 2.0 & 1176.9 & 7.722 \\
Refined time step & 1.0 & 1169.3 & 7.724 \\
Refined mesh (band) & 2.0 & 1173.2 & 8.819 \\
Refined mesh + time step & 1.0 & 1174.9 & 8.819 \\
\bottomrule
\end{tabular}
\end{table}

\noindent Halving the time step at the baseline mesh changes $|w|_{\max}$ by only \SI{0.03}{\%} and $T_{\mathrm{peak}}$ by \SI{0.65}{\%}. Refining the heating-band mesh from $h_{\mathrm{ref}}=\SI{5}{mm}$ to \SI{2.5}{mm} increases $|w|_{\max}$ by \SI{14.2}{\%} while leaving $T_{\mathrm{peak}}$ within \SI{0.3}{\%} of the baseline. The combined refined mesh and time-step run differs from the refined-mesh-only result by less than \SI{0.01}{\%} in $|w|_{\max}$, so the present sensitivity matrix indicates that local mesh resolution near the heat path is the dominant discretization factor for this case.

\begin{table}[H]
\centering
\caption{Diagnostic comparison between the calibrated inherent-strain surrogate and the optional J2 elastoplastic mechanics path for representative Li2023 Case~6 ($v=\SI{11}{mm/s}$, $E=\SI{600}{J/mm}$, $t=\SI{14}{mm}$, two passes).}
\label{tab:li2023-elastoplastic-spotcheck}
\small
\resizebox{\linewidth}{!}{%
\begin{tabular}{lrrl}
\toprule
Model & $|w|_{\max}$ (mm) & $T_{\mathrm{peak}}$ (K) & Plasticity indicator \\
\midrule
Inherent strain (calibrated) & 7.722 & 1176.9 & $\varepsilon_0=4.7736\times 10^{-2}$ surrogate amplitude; no resolved plastic-strain field \\
Elastoplastic (J2, strong-coupled) & 0.096 & 1176.9 & $\bar{\varepsilon}^{p}_{\max}=5.8846\times 10^{-2}$; mean $4.3067\times 10^{-4}$; active elements 5342/36780 (14.5\%) \\
\bottomrule
\end{tabular}
}
\end{table}

\noindent The J2 run advances equivalent plastic strain through the saved thermal history using the same mesh and thermal loading as the calibrated Case~6 run. It confirms that the thermal cycle can trigger localized yielding, but the present optional elastoplastic path does not reproduce the calibrated residual-deflection magnitude. The comparison is therefore retained as a diagnostic verification item rather than treated as independent validation of the inherent-strain calibration.

\begin{table}[H]
\centering
\caption{Summary validation metrics. MAE, RMSE, and bias are in mm; MAPE and maximum error are percentages.}
\label{tab:validation-metrics}
\small
\resizebox{\linewidth}{!}{%
\begin{tabular}{lrrrrrr}
\toprule
Run / split & $n$ & MAE & RMSE & Bias & MAPE & Max $|\%|$ \\
\midrule
v9 moving torch, all cases & 11 & 0.077 & 0.094 & 0.073 & 1.17 & 2.78 \\
v9 moving torch, calibration anchors & 4 & 0.034 & 0.053 & 0.022 & 0.52 & 1.77 \\
v9 moving torch, non-anchor validation & 7 & 0.102 & 0.112 & 0.102 & 1.55 & 2.78 \\
v10 line source, collapsed + quench & 11 & 2.578 & 3.862 & $-$2.468 & 23.56 & 50.13 \\
v11 line source, sequential + quench & 5 & 0.034 & 0.045 & 0.020 & 0.27 & 0.66 \\
\bottomrule
\end{tabular}
}
\end{table}

\Cref{fig:li2023-parity,fig:li2023-residuals,fig:li2023-speed-line} visualize the moving-torch validation. Calibration anchors are marked separately from non-anchor validation cases to make the fitted cases explicit.

\begin{figure}[H]
\centering
\includegraphics[width=0.72\linewidth]{li2023_v9_parity.png}
\caption{Parity plot for the v9 moving-torch baseline. Calibration anchors are shown separately from non-anchor validation cases.}
\label{fig:li2023-parity}
\end{figure}

\begin{figure}[H]
\centering
\includegraphics[width=0.78\linewidth]{li2023_v9_residuals.png}
\caption{Case-wise residual error for the v9 moving-torch baseline. The largest residual occurs in the low-energy single-pass speed sweep.}
\label{fig:li2023-residuals}
\end{figure}

\begin{figure}[H]
\centering
\includegraphics[width=0.72\linewidth]{li2023_speed_sweep_v9.png}
\caption{Speed-sweep response for cases 1--5. Experiments show a small decrease in deflection with speed; the current inherent-strain law predicts an almost flat response.}
\label{fig:li2023-speed-line}
\end{figure}

\subsection{Full-Line Thermal-Flux and Quench Comparison}
\label{sec:validation-induction}

Two full-line thermal-flux variants were run to test whether a reduced line-source approximation can reproduce the moving-torch forming response. The first applies a single full-line heating event with quenching and collapses an $N$-pass plan into one longer event using an effective velocity $v_{\mathrm{eff}}=v/N$. The second repeats the same full-line heating event sequentially for the specified pass count while retaining quenching. The sequential variant therefore preserves the distinct thermal cycles present in the Li2023 multi-pass experiments. This comparison evaluates process scheduling and heat-input representation only; it is not an electromagnetic validation study.

\begin{table}[H]
\centering
\caption{Process comparison for Li2023 cases: moving torch baseline (v9), collapsed full-line thermal flux with quench (v10), and sequential full-line thermal flux with quench (v11). Values are predicted residual deflection in mm with percentage error in parentheses.}
\label{tab:li2023-process}
\scriptsize
\resizebox{\textwidth}{!}{%
\begin{tabular}{rrrrrrrrr}
\toprule
Case & $v$ & $E$ & $t$ & $N$ & $w_{\mathrm{exp}}$ & v9 torch & v10 line+quench & v11 line+quench \\
 & (mm/s) & (J/mm) & (mm) & & (mm) & (mm) & (mm) & (mm) \\
\midrule
1  & 6  & 600  & 14 & 1 & 5.785 &  5.888 (+1.8\%) &  5.879 (+1.6\%) & -- \\
2  & 7  & 600  & 14 & 1 & 5.771 &  5.888 (+2.0\%) &  5.879 (+1.9\%) & -- \\
3  & 8  & 600  & 14 & 1 & 5.760 &  5.887 (+2.2\%) &  5.879 (+2.1\%) & -- \\
4  & 9  & 600  & 14 & 1 & 5.748 &  5.886 (+2.4\%) &  5.879 (+2.3\%) & -- \\
5  & 10 & 600  & 14 & 1 & 5.727 &  5.886 (+2.8\%) &  5.879 (+2.7\%) & -- \\
6  & 11 & 600  & 14 & 2 & 7.712 &  7.722 (+0.1\%) &  3.898 ($-$49.5\%) &  7.713 (+0.0\%) \\
7  & 12 & 600  & 14 & 2 & 7.696 &  7.722 (+0.3\%) &  3.898 ($-$49.4\%) &  7.713 (+0.2\%) \\
87 & 12 & 1000 & 10 & 2 & 13.506 & 13.486 ($-$0.2\%) &  6.735 ($-$50.1\%) & 13.471 ($-$0.3\%) \\
88 & 13 & 1000 & 10 & 2 & 13.443 & 13.486 (+0.3\%) &  6.735 ($-$49.9\%) & 13.471 (+0.2\%) \\
89 & 14 & 1000 & 10 & 2 & 13.383 & 13.485 (+0.8\%) &  6.735 ($-$49.7\%) & 13.471 (+0.7\%) \\
90 & 15 & 1000 & 10 & 1 & 11.322 & 11.319 ($-$0.0\%) & 11.310 ($-$0.1\%) & -- \\
\bottomrule
\end{tabular}
}
\end{table}

\noindent The collapsed full-line variant (v10) matches single-pass cases but underpredicts multi-pass cases by about \SI{50}{\%}. This is a process-scheduling error: one longer full-line heat does not reproduce two separate thermal cycles and associated inherent-strain accumulation. The sequential full-line variant (v11) corrects this for the multi-pass subset (cases 6, 7, 87--89), with MAPE \SI{0.27}{\%}. These results support the reduced full-line thermal source as a useful process-planning approximation only when the physical pass schedule is preserved.

\begin{figure}[H]
\centering
\includegraphics[width=0.86\linewidth]{li2023_process_comparison.png}
\caption{Process comparison for multi-pass cases. The collapsed line-source schedule (v10) gives roughly half the measured deflection; preserving sequential passes (v11) recovers the moving-torch and experimental response.}
\label{fig:li2023-process-comparison}
\end{figure}

\subsection{Keel-Plate Demonstration Cases}
\Cref{tab:keel} summarizes keel-plate demonstration runs (summarized from the per-run output files in the project repository). These are application demonstrations rather than experimental validation cases. Runs with physically inconsistent peak temperatures from earlier unit-normalization tests are retained in the repository but excluded from the main quantitative table.

\begin{table}[H]
\centering
\caption{Keel-plate demonstration cases retained for publication discussion: heat mode, passes, peak flux $q_0$, radius $r_0$, peak temperature, and deflection extrema.}
\label{tab:keel}
\scriptsize
\setlength{\tabcolsep}{4pt}
\renewcommand{\arraystretch}{1.05}
\resizebox{\linewidth}{!}{%
\begin{tabular}{llrrrrl}
\toprule
Case & Heat mode & Passes & $q_0$ (W/mm$^2$) & $r_0$ (mm) & $T_{\max}$ (K) & $w_{\max}/w_{\min}$ (mm) \\
\midrule
Keel\_fast\_parallel & simultaneous & 1 & 59.05 & 9.0 & 1173.0 & 0.003/$-0.044$ \\
Keel\_practical      & simultaneous & 1 & 96.99 & 9.0 & 1171.7 & 8.32/$-1.95$ \\
Keel\_li2023\_scaled & sequential   & 2 & 55.42 & 7.0 & 1476.0 & 20.8/$-7.96$ \\
\bottomrule
\end{tabular}%
}
\end{table}

The demonstration deflection ranges are \SI{0.047}{mm}, \SI{10.27}{mm}, and \SI{28.79}{mm} for the fast-parallel, practical, and Li2023-scaled keel cases, respectively. The keel examples demonstrate the workflow on larger geometry; publication-grade keel validation requires a measured or target curvature profile and a completed thermal-unit audit.
